\section{Least preimages}\label{sec:companion}

We now prove Theorem~\ref{thm:companion}.  The prime that dominates a
basic representation also determines the interval in which its least
preimage lies.  Counting that prime first gives the weight \(f_k\).
The resulting approximation is additive; a separate argument will show
when its coefficient is bounded away from zero.

Throughout this section, \(k\geq1\) is a fixed integer.  The mass
\(M_f(x;H)\) is defined in \eqref{eq:mass}; its summands correspond to
distinct data \((d,p_1,\ldots,p_R)\), regardless of how many basic
witnesses yield the same data.  Write
\[
 f_k(r)=
 \begin{cases}
  0,&1\leq r\leq k,\\
  1-k/r,&k<r<k+1,\\
  1/r,&r\geq k+1.
 \end{cases}
\]
Thus \(0\leq f_k\leq1/(k+1)\), and \(f_k(r)>0\) exactly when \(r>k\).
Equivalently, \(f_k(r)\) is the length of
\[
 \{t\in[0,1]:k<rt\leq k+1\}.
\]
It is therefore the permitted fraction of the normalized value interval
when the least-preimage ratio is fixed at \(r\).

\subsection{Counting the interval for the largest prime}

\begin{lemma}\label{lem:weighted-endpoint}
For each sufficiently large fixed \(H\), as \(x\) tends to infinity,
\begin{equation}\label{eq:weighted-mass}
 \left|N_k(x)-\frac{x}{\log x}M_{f_k}(x;H)\right|
 \leq\bigl(\epsilon_{H,k}+o_{x\to\infty;H,k}(1)\bigr)
       \frac{xG_m}{\log x},
\end{equation}
where \(\epsilon_{H,k}\geq0\) tends to zero as \(H\to\infty\).
\end{lemma}

\begin{proof}
Proposition~\ref{prop:unique-prefix} allows us, with an error of
\((\epsilon_H+o_{x\to\infty;H}(1))xG_m/\log x\), to count tuples
\((p_0,\ldots,p_R,d)\) satisfying
\[
 (p_0-1)\cdots(p_R-1)d\leq x,
 \qquad
 kx<p_0\cdots p_R\ell(d)\leq(k+1)x.
\]
Both the exceptional values and the tuples above them are controlled
there.  This use of the least-preimage identity requires that every
preimage have the same prefix.  The identity would not follow from the
existence of a single basic witness.

Fix the remaining data \((d,p_1,\ldots,p_R)\), and put
\[
 D=d\prod_{i=1}^R(p_i-1),\qquad
 A=\ell(d)\prod_{i=1}^Rp_i,\qquad
 r=\frac{\ell(d)}d,\qquad
 q=\prod_{i=1}^R\left(1-\frac1{p_i}\right).
\]
Here \(r\geq1\).  Set \(T=x/D\) and \(U=x/A=Tq/r\).
The admissible largest primes satisfy exactly
\begin{equation}\label{eq:largest-prime-interval}
 p_0\geq x^{0.9},\qquad
 kU<p_0\leq\min\{T+1,(k+1)U\}.
\end{equation}
Once the lower bound for \(p_0\) holds, the basic-witness restrictions
on the remaining data do not change with \(p_0\); recall that their
zeroth simplex coordinate is \(B\).

Ignoring the lower cutoff and replacing \(T+1\) by \(T\), the interval
has length
\[
 \min\{T,(k+1)U\}-\min\{T,kU\}=T f_k(r/q).
\]
The replacement changes its length by at most one.  Moreover,
\(z\mapsto f_k(1/z)\) is \((k+1)\)-Lipschitz on \([0,1]\),
where its value at zero is zero.  Consequently
\begin{equation}\label{eq:weight-perturbation}
 |f_k(r/q)-f_k(r)|
 \leq\frac{(k+1)(1-q)}r
 \leq(k+1)\sum_{i=1}^R\frac1{p_i}.
\end{equation}

For fixed \(H\), the deterministic size bounds in
\eqref{eq:cofactor-size} give \(T=x^{1-o(1)}\) uniformly in the
remaining data.  The tail has a witness \(w\) bounded in terms of
\(H\) only, so \(\ell(d)\leq w\) bounds \(r\) in terms of \(H\).
The prime bands bound \(q\) away from zero.  Hence
\(U=x^{1-o(1)}\) uniformly as well, and \(kU>x^{0.9}\) for all
sufficiently large \(x\).  The lower cutoff in
\eqref{eq:largest-prime-interval} is then inactive.

Apply the prime number theorem at its two endpoints, taking the positive
part of the difference if the interval is empty.  Each nonzero endpoint
is \(x^{1-o(1)}\), and each is \(O_k(T)\).  Thus this gives, uniformly
in the remaining data,
\[
 \#\{p_0\text{ satisfying \eqref{eq:largest-prime-interval}}\}
 =\frac{T}{\log x}
   \left(f_k(r)+O_k\left(\sum_{i=1}^R\frac1{p_i}\right)
                +o_{x\to\infty;H,k}(1)\right).
\]
The shift by one is included in the last error.  This is an additive
estimate relative to \(T/\log x\), obtained from the ordinary prime
number theorem; no relative estimate on a short prime interval is
required.

The geometric separation of the bands implies
\[
 \sup\sum_{i=1}^R\frac1{p_i}
       \ll \exp(-cH\rho^{-H}).
\]
Indeed \(b_{i-1}/b_i\geq\rho^{-1}\), and
\(p_i\geq\exp\exp(c b_i)\); the resulting reciprocal series is
bounded by a constant times its smallest-scale bound.  In particular,
this estimate is uniform in the growing number \(R\) of primes.
Since \(M_1\ll K_HG_m\) and \(K_H=H^{o(1)}\), summing the displayed
prime count gives an error bounded by
\[
 \left(O_k\bigl(K_H\exp(-cH\rho^{-H})\bigr)
           +o_{x\to\infty;H,k}(1)\right)\frac{xG_m}{\log x}.
\]
Together with the tuple-to-value error this proves
\eqref{eq:weighted-mass}.
\end{proof}

The additive form is necessary at this stage.  For example, when
\(r=k\) and \(q<1\), one has \(f_k(r)=0\) but \(f_k(r/q)>0\).
Inequality~\eqref{eq:weight-perturbation} controls this boundary without
asserting a relative approximation to a zero term.

Apply Lemma~\ref{lem:coefficient-convergence} with the function
\[
 C_{f_k}(x)=\frac{N_k(x)\log x}{xG_m},
\]
which is independent of \(H\).  Lemma~\ref{lem:weighted-endpoint}
provides its required comparison with \(M_{f_k}/G_m\).  We obtain the
uniform convergence of \(A_H(f_k;s)\) on \(0\leq s<1\) and
\begin{equation}\label{eq:companion-additive}
 N_k(x)=\frac{xG_m}{\log x}
          \bigl(A(f_k;\theta)+o(1)\bigr).
\end{equation}
It remains to distinguish a coefficient bounded away from zero from an
identically zero count.

\subsection{Seed propagation and bounded ratios}
The positivity argument needs two distinct facts.  A seed with
\(\ell(d)/d>k\) must produce a positive proportion of values whose
ratios stay away from \(k\).  These ratios must also remain bounded
on a positive proportion of values, since \(f_k(r)\to0\) as
\(r\to\infty\).  The next two lemmas supply these facts.

Preservation of a complete inverse fiber under multiplication by a prime
already appears in Erd\H{o}s's proof of Theorem~4 \cite{Erdos1958}.
Ford's construction gives such preservation on a positive proportion
of the distinct-totient scale. A later simultaneous two-seed form appears
in Pollack, Pomerance and Trevi\~no \cite[Lemma~4.1]{PPT2013},
again using Ford's construction. We use the following single-seed
consequence of Ford's proof.


\begin{lemma}[Ford's inverse-fiber propagation]\label{lem:fiber-propagation}
Fix a totient \(d\), and write
\(\varphi^{-1}(d)=\{d_1,\ldots,d_\kappa\}\).
There are constants \(\eta_d>0\) and \(x_d\) such that, for every
\(x\geq x_d\), at least \(\eta_dV(x)\) distinct totients \(v\leq x\)
have the form
\[
 v=d\varphi(b),\qquad
 \varphi^{-1}(v)=\{bd_1,\ldots,bd_\kappa\}
\]
for some positive integer \(b\).  In particular, these values satisfy
\(\ell(v)/v\geq\ell(d)/d\).
\end{lemma}

\begin{proof}
Use the construction in the proof of Theorem~2 of
\cite[Section~5, pp.~24--26, Equations~(5.10)--(5.19)]{FordTotients2013}.
For the fixed seed \(d\), Ford constructs a set \(\mathcal B\) and
excludes at most half its members.  For each remaining \(b\), every
preimage of \(d\varphi(b)\) is one of the products \(bd_i\).
The paragraph following Equation~(5.18) also proves that these members
give distinct totients.  The concluding bound on p.~26 is
\(\lvert\mathcal B\rvert/2\gg_\varepsilon d^{-1-\varepsilon}V(x)\)
for sufficiently large \(x\).  This full-fiber conclusion is recorded
explicitly in \cite[Section~7.3, p.~39]{FordTotients2013}.
Finally,
\[
 \frac{\ell(v)}v
    =\frac{b}{\varphi(b)}\frac{\ell(d)}d
    \geq\frac{\ell(d)}d.
\]
\end{proof}

We also need to keep these ratios in a bounded interval.  This follows
from the existence of a bounded arithmetic tail and does not require
uniqueness of the prefix.

\begin{lemma}\label{lem:ratio-tightness}
For every \(\varepsilon>0\), there are \(C<\infty\) and \(x_0\)
such that
\[
 \#\{v\in\mathcal V:v\leq x,\ \ell(v)/v>C\}
     \leq\varepsilon V(x)\qquad(x\geq x_0).
\]
\end{lemma}

\begin{proof}
Choose one sufficiently large fixed cut \(H_0\) in
Proposition~\ref{prop:basic}, so that its exceptional set has at most
\(\varepsilon V(x)\) members for all sufficiently large \(x\).
Every remaining value has a basic witness
\(n=p_0\cdots p_{R_0}w\), where \(R_0=m-H_0\),
\(w<p_{R_0}\), and \(w\) is bounded in terms of \(H_0\) only.
The prefix primes are distinct, and their bands give a uniform bound
for \(\prod_{i=0}^{R_0}p_i/(p_i-1)\).  Therefore
\[
 \frac{\ell(v)}v\leq\frac{n}{\varphi(n)}
   =\frac{w}{\varphi(w)}
      \prod_{i=0}^{R_0}\frac{p_i}{p_i-1}\leq C(H_0).
\]
The last bound is independent of \(x\), as required.
\end{proof}

\subsection{The positive and zero alternatives}

Suppose first that there is a totient \(d_*\) with
\(\ell(d_*)>kd_*\).  Choose \(\delta>0\) such that
\(\ell(d_*)/d_*>k+2\delta\).
Lemma~\ref{lem:fiber-propagation} gives a fixed \(\eta>0\) such that
at least \(\eta V(x)\) values satisfy
\(\ell(v)/v>k+2\delta\) for every sufficiently large \(x\).
Apply Lemma~\ref{lem:ratio-tightness} with \(\varepsilon=\eta/4\).
It supplies a fixed \(C\) for which at least
\(3\eta V(x)/4\) values satisfy
\begin{equation}\label{eq:bounded-positive-ratios}
 k+2\delta<\frac{\ell(v)}v\leq C.
\end{equation}
The numbers \(\eta,\delta,C\) are fixed before the cut \(H\) is
allowed to grow.

For every sufficiently large \(H\), and then sufficiently large \(x\),
Proposition~\ref{prop:unique-prefix} removes at most \(\eta V(x)/4\)
of these values.  For each remaining value its unique tuple satisfies
\[
 \frac{\ell(v)}v=\frac{\ell(d)}d
                 \prod_{i=0}^R\frac{p_i}{p_i-1}.
\]
The product is \(1+O(\epsilon_H)+o_{x\to\infty;H}(1)\), uniformly
in the tuple.  Increasing \(H\) and then \(x\), we may bound it by
\((k+2\delta)/(k+\delta)\).  Thus at least \(\eta V(x)/2\)
distinct tuples have their tail ratio in the fixed interval
\[
 I=[k+\delta,C].
\]

Let \(M_{\one_I}\) denote the mass \eqref{eq:mass} with weight
\(\one_I(r)\).  The ordinary prime number theorem, counting all
admissible largest primes for data whose tail ratio lies in \(I\),
bounds the number of such tuples above by
\[
 \bigl(1+o_{x\to\infty;H}(1)\bigr)
       \frac{x}{\log x}M_{\one_I}(x;H).
\]
Comparison with \(\eta V(x)/2\), followed by
\eqref{eq:ford-scale}, gives
\[
 M_{\one_I}(x;H)\geq c_1G_m
\]
for all sufficiently large \(H\) and then sufficiently large \(x\),
where \(c_1>0\) is independent of \(H\).  Since
\[
 \inf_{r\in I}f_k(r)
      \geq\min\left\{\frac{\delta}{k+\delta},\frac1C\right\}>0,
\]
we conclude that
\begin{equation}\label{eq:weighted-mass-positive}
 M_{f_k}(x;H)\geq c_2G_m
\end{equation}
with \(c_2>0\) independent of sufficiently large \(H\).

We transfer this lower bound first to the normalized count
\(C_{f_k}(x)=N_k(x)\log x/(xG_m)\), which does not depend on \(H\).
Fix a sufficiently large \(H\) such that the error
\(\epsilon_{H,k}\) in \eqref{eq:weighted-mass} is less than \(c_2/4\),
and then take \(x\) large enough that its remaining error is less
than \(c_2/4\).  Equations~\eqref{eq:weighted-mass} and
\eqref{eq:weighted-mass-positive} give
\(C_{f_k}(x)\geq c_2/2\) for every sufficiently large \(x\).
This count is also bounded above, since \(N_k(x)\leq V(x)\) and
\eqref{eq:ford-scale} holds.  The bounds clause of
Lemma~\ref{lem:coefficient-convergence} therefore gives
\[
 \inf_{0\leq s<1}A(f_k;s)\geq c_2/2>0.
\]
Therefore the additive formula \eqref{eq:companion-additive} is
multiplicative in this case.  Its positive lower bound, the inequality
\(N_k(x)\leq V(x)\), and \eqref{eq:ford-scale} give
\(N_k(x)\asymp_k V(x)\).

If instead \(\ell(d)\leq kd\) for every totient \(d\), then
\(\ell(v)\leq kv\leq kx\) for every \(v\leq x\).  Hence
\(N_k(x)=0\) for every \(x\).  Every weight
\(f_k(\ell(d)/d)\) in the finite arithmetic formula is also zero,
so \(A_H(f_k;s)=A(f_k;s)=0\).  This proves the dichotomy in Theorem~\ref{thm:companion}.
Its final assertion is established next.

\begin{remark}\label{rem:least-preimage-scope}
The positive alternative occurs for \(k=1\) and \(k=2\).  Ford gives
\(d=2^{18}\cdot257\) as a totient all of whose preimages are divisible
by \(8\) \cite[Section~7.3, p.~39]{FordTotients2013}.  Each such
preimage \(n\) is even and has an odd prime factor: a power of \(2\)
would have a power of \(2\) as its totient.  Thus
\(n/\varphi(n)=\prod_{p\mid n}p/(p-1)>2\), and \(\ell(d)>2d\).

We do not determine all integers \(k\) for which the positive
alternative occurs.  Its occurrence for every integer \(k\geq1\)
is equivalent to the unboundedness of \(\ell(d)/d\) over totients
\(d\).  Neither that unboundedness nor the complete classification
of \(k\) is established here.
\end{remark}
